
\documentclass[11pt]{article}
\usepackage[margin=0.8in]{geometry}
\usepackage{amsmath,amssymb,booktabs,array,longtable}
\usepackage{hyperref}
\usepackage{enumitem}
\setlist{nosep}

\begin{document}

\begin{center}
{\Large\textbf{Shiv Nadar University Chennai}}\\
\textbf{CS3807 -- Deep Learning Laboratory}\\
\textbf{Experiment 2}\\
{\large\textbf{Implementation of a Multi-Layer Perceptron (MLP) for Multi-Class Image Classification}}
\end{center}

\renewcommand{\arraystretch}{1.25}

\begin{tabular}{|p{3.8cm}|p{8cm}|p{2cm}|}
\hline
Degree \& Branch & B.Tech Artificial Intelligence \& Data Science & Semester V\\ \hline
Subject Code \& Name & CS3807 -- Deep Learning Laboratory & AY: 2026--27\\ \hline
\end{tabular}

\section*{1. Objective}
Implement an MLP using TensorFlow/Keras on the Fashion-MNIST dataset. Learn image preprocessing, flattening, model construction, training, evaluation, and automated hyperparameter optimization.

\section*{2. Learning Outcomes}
Students will be able to:
\begin{enumerate}
\item Explain the architecture of an MLP.
\item Preprocess image datasets.
\item Build and train an MLP.
\item Evaluate classification performance.
\item Perform automated hyperparameter optimization.
\item Recommend the best model configuration.
\end{enumerate}

\section*{3. Background Theory}
An MLP consists of an input layer, hidden layers and an output layer.

\[
z^{(l)}=W^{(l)}a^{(l-1)}+b^{(l)},\qquad
a^{(l)}=f(z^{(l)})
\]

Output layer:
\[
\hat y=\mathrm{Softmax}(z)
\]

Loss:
\[
L=-\sum_i y_i\log(\hat y_i)
\]

Recommended architecture:

\begin{center}
784 $\rightarrow$ Dense(128, ReLU) $\rightarrow$ Dense(64, ReLU) $\rightarrow$ Dense(10, Softmax)
\end{center}

\section*{4. Dataset}
Dataset: Fashion-MNIST

Training Images: 60,000 \\
Testing Images: 10,000 \\
Classes: 10 \\
Image Size: $28\times28$

\section*{5. Image Preprocessing}

Fashion-MNIST images are stored as $28\times28$ matrices. Dense layers require a one-dimensional feature vector.

\[
28\times28=784
\]

Example:

\[
\begin{bmatrix}
10&20\\30&40
\end{bmatrix}
\rightarrow
[10,20,30,40]
\]

Perform the following preprocessing:

\begin{enumerate}
\item Load Fashion-MNIST.
\item Display sample images.
\item Flatten images.
\item Normalize pixels to [0,1].
\item Convert labels into one-hot vectors.
\end{enumerate}

\section*{6. Experimental Procedure}

\textbf{Task 1: Dataset Exploration}
\begin{itemize}
\item Load Fashion-MNIST.
\item Print dataset dimensions.
\item Display ten sample images.
\item Plot class distribution.
\end{itemize}

\textbf{Expected Output:}
Dataset summary and sample images.

\bigskip

\textbf{Task 2: Data Preprocessing}

\begin{itemize}
\item Flatten images.
\item Normalize pixel values.
\item Print tensor shapes before and after preprocessing.
\end{itemize}

\bigskip

\textbf{Task 3: Model Construction}

Construct an MLP with

\begin{itemize}
\item Input layer (784)
\item Dense(128, ReLU)
\item Dense(64, ReLU)
\item Dense(10, Softmax)
\end{itemize}

Display \texttt{model.summary()}.

\bigskip

\textbf{Task 4: Model Training}

Compile using

\begin{itemize}
\item Optimizer: Adam
\item Loss: Categorical Cross Entropy
\item Metric: Accuracy
\end{itemize}

Train for 20 epochs with batch size 32.

\bigskip

\textbf{Task 5: Model Evaluation}

Compute

\begin{itemize}
\item Accuracy
\item Precision
\item Recall
\item F1-score
\item Confusion Matrix
\item Classification Report
\end{itemize}

\section*{7. Hyperparameter Optimization}

Instead of manually selecting hyperparameters, students shall perform automated hyperparameter optimization using either \textbf{GridSearchCV} or \textbf{RandomizedSearchCV} (recommended) together with the SciKeras wrapper.

Optimize the following search space.

\begin{center}
\begin{tabular}{|p{5cm}|p{7cm}|}
\hline
Hyperparameter & Candidate Values\\
\hline
Hidden Layers & 1, 2, 3\\
Hidden Neurons & 32, 64, 128, 256\\
Learning Rate & 0.1, 0.01, 0.001\\
Batch Size & 16, 32, 64, 128\\
Epochs & 10, 20, 30\\
Optimizer & SGD, Adam, RMSProp\\
Activation Function & ReLU, Tanh, Sigmoid\\
Dropout Rate (Optional) & 0.0, 0.2, 0.5\\
\hline
\end{tabular}
\end{center}

\subsection*{Tasks}

\begin{enumerate}
\item Build a baseline MLP model.
\item Define the hyperparameter search space.
\item Perform Grid Search or Randomized Search using 5-fold cross-validation.
\item Record the best hyperparameter combination.
\item Retrain the model using the optimized hyperparameters.
\item Evaluate the optimized model on the testing dataset.
\item Compare the optimized model with the baseline model.
\end{enumerate}

\section*{8. Mandatory Plots}

Generate the following plots.

\begin{enumerate}
\item Sample Images
\item Class Distribution
\item Training Accuracy vs Epoch
\item Validation Accuracy vs Epoch
\item Training Loss vs Epoch
\item Validation Loss vs Epoch
\item Confusion Matrix
\item Hyperparameter Search Results
\item Best Model Accuracy Comparison
\end{enumerate}

Write a 2--3 line inference for each plot.

\section*{9. Results}

\textbf{Best Hyperparameters}

\begin{tabular}{|p{6cm}|p{6cm}|}
\hline
Hidden Layers &\\\hline
Hidden Neurons &\\\hline
Learning Rate &\\\hline
Batch Size &\\\hline
Optimizer &\\\hline
Activation Function &\\\hline
Epochs &\\\hline
Dropout &\\\hline
Cross-validation Accuracy &\\\hline
Testing Accuracy &\\\hline
\end{tabular}

\vspace{3mm}

\textbf{Performance Comparison}

\begin{tabular}{|p{4cm}|c|c|}
\hline
Metric & Baseline & Optimized\\
\hline
Accuracy&&\\
Precision&&\\
Recall&&\\
F1-score&&\\
Training Time&&\\
\hline
\end{tabular}

\section*{10. Discussion}

Answer the following:

\begin{enumerate}
\item Which optimization method was used?
\item Which hyperparameters were selected?
\item Did optimization improve performance?
\item Which hyperparameter had the greatest impact?
\item Compare Grid Search and Randomized Search.
\item Recommend the best MLP configuration.
\end{enumerate}

\section*{11. Report Contents}

Include Objective, Theory, Dataset, Experimental Procedure, Source Code, Hyperparameter Optimization, Results, Plots with Inference, Discussion, Conclusion and References.

\section*{12. References}
\begin{enumerate}
\item Goodfellow et al., \emph{Deep Learning}.
\item Bishop, \emph{Pattern Recognition and Machine Learning}.
\item Haykin, \emph{Neural Networks and Learning Machines}.
\item Fashion-MNIST Dataset.
\item TensorFlow/Keras Documentation.
\end{enumerate}

\end{document}
